% SPDX-License-Identifier: MPL-2.0
% SPDX-FileCopyrightText: 2026 Jonathan D.A. Jewell
%
% Consent-Aware Programming: Embedding Ethical Deliberation, Attribution,
% and Emotional Context in Language Semantics
%
% Target: arXiv preprint (cs.PL / cs.HC cross-list)

\documentclass[11pt,a4paper]{article}

% --- Packages ---
\usepackage[utf8]{inputenc}
\usepackage[T1]{fontenc}
\usepackage{amsmath,amssymb,amsthm}
\usepackage{mathpartir}
\usepackage{stmaryrd}
\usepackage{listings}
\usepackage{xcolor}
\usepackage{hyperref}
\usepackage{cleveref}
\usepackage{booktabs}
\usepackage{graphicx}
\usepackage{natbib}
\usepackage{microtype}
\usepackage{enumitem}
\usepackage{tikz}
\usetikzlibrary{arrows.meta,positioning,shapes.geometric,calc}
\usepackage{fancyvrb}
\usepackage[margin=1in]{geometry}

% --- Theorem environments ---
\newtheorem{theorem}{Theorem}[section]
\newtheorem{lemma}[theorem]{Lemma}
\newtheorem{definition}[theorem]{Definition}
\newtheorem{proposition}[theorem]{Proposition}
\newtheorem{corollary}[theorem]{Corollary}
\newtheorem{example}[theorem]{Example}

% --- Listings configuration ---
\definecolor{wokekw}{RGB}{120,50,150}
\definecolor{wokestr}{RGB}{0,120,60}
\definecolor{wokecmt}{RGB}{100,100,100}
\definecolor{wokeann}{RGB}{180,60,30}
\definecolor{codebg}{RGB}{248,248,248}

\lstdefinelanguage{WokeLang}{
  morekeywords={to,remember,only,if,okay,thanks,measured,in,data,then,
                clean,analyze,feeling,confident,return,true,false,
                when,otherwise,for,each,while,match,with,import,from,
                module,expose,type,is,has,end,pipeline,step,consent,
                revoke,grant,audit,acknowledge},
  sensitive=true,
  morecomment=[l]{//},
  morecomment=[s]{/*}{*/},
  morestring=[b]",
  moredelim=[s][\color{wokeann}]{@}{)},
}

\lstset{
  language=WokeLang,
  basicstyle=\small\ttfamily,
  keywordstyle=\color{wokekw}\bfseries,
  stringstyle=\color{wokestr},
  commentstyle=\color{wokecmt}\itshape,
  backgroundcolor=\color{codebg},
  frame=single,
  framerule=0.4pt,
  rulecolor=\color{gray!40},
  numbers=left,
  numberstyle=\tiny\color{gray},
  numbersep=8pt,
  breaklines=true,
  breakatwhitespace=true,
  tabsize=2,
  showstringspaces=false,
  captionpos=b,
  xleftmargin=1.5em,
  framexleftmargin=1.5em,
  aboveskip=1em,
  belowskip=0.8em,
}

% --- Semantic notation ---
\newcommand{\sem}[1]{\llbracket #1 \rrbracket}
\newcommand{\eval}{\Downarrow}
\newcommand{\step}{\longrightarrow}
\newcommand{\consent}{\mathsf{consent}}
\newcommand{\granted}{\mathsf{granted}}
\newcommand{\denied}{\mathsf{denied}}
\newcommand{\pending}{\mathsf{pending}}
\newcommand{\unitty}[1]{\langle #1 \rangle}
\newcommand{\feeling}[2]{\mathsf{feel}(#1, #2)}
\newcommand{\attr}[2]{\mathsf{attr}(#1, #2)}
\newcommand{\WokeLang}{\textsc{WokeLang}}

% --- Title ---
\title{%
  \textbf{Consent-Aware Programming:}\\[0.3em]
  \large Embedding Ethical Deliberation, Attribution, and Emotional Context\\
  in Language Semantics%
}

\author{%
  Jonathan D.A.\ Jewell\\
  \texttt{j.d.a.jewell@open.ac.uk}\\
  \texttt{https://github.com/hyperpolymath}%
}

\date{March 2026}

\begin{document}

\maketitle

% ============================================================================
% ABSTRACT
% ============================================================================
\begin{abstract}
Programming languages encode assumptions about what matters.
Most languages assume that what matters is \emph{computation}---the
transformation of inputs to outputs, the allocation and reclamation of
memory, the dispatch of messages between processes.  We argue that this
assumption is incomplete.  Programs are social artifacts: they are
written by people, for people, within institutional and ethical
contexts that shape their meaning no less than their operational
semantics.

We present \WokeLang{}, a programming language that embeds ethical
deliberation, attribution, and emotional context directly into its
semantics.  \WokeLang{} introduces four novel constructs:
\emph{consent gates} (\texttt{only if okay}), which create mandatory
deliberation points around ethically significant operations;
\emph{units of measure} (\texttt{measured in}), which carry
dimensional information as first-class types oriented toward human
communication; \emph{gratitude blocks} (\texttt{thanks to}), which
make contribution attribution an executable, auditable program
construct; and \emph{empathy annotations} (\texttt{@feeling}), which
allow developers to attach confidence and emotional metadata to code.
We provide formal operational semantics for consent gates, describe
the integration of these features into a Hindley-Milner type system
with dimensional analysis, and report on an implementation targeting
WebAssembly via an OCaml frontend and Rust runtime.

Our design philosophy holds that programming languages are
\emph{communication media} between humans, not merely instruction sets
for machines.  We draw on values-sensitive design, feminist HCI, and
the philosophy of technology to argue that language designers have both
the opportunity and the responsibility to make ethical deliberation
syntactically cheap and semantically meaningful.
\end{abstract}

\vspace{0.5em}
\noindent\textbf{Keywords:} programming language design, consent,
ethics in computing, human-centered programming, units of measure,
attribution, values-sensitive design, formal semantics

% ============================================================================
% 1. INTRODUCTION
% ============================================================================
\section{Introduction}
\label{sec:introduction}

Every programming language embodies a theory of what programming
\emph{is}.  C treats programming as the management of memory and
machine words.  Haskell treats it as the composition of mathematical
functions.  Erlang treats it as the coordination of isolated,
message-passing processes.  These are not merely implementation
choices; they are \emph{ontological commitments} that shape how
programmers think about the systems they build
\citep{sapir1929status,whorf1956language,iverson1980notation}.

What no mainstream language treats programming as is a fundamentally
\emph{social} and \emph{ethical} activity.  Yet programs are written by
teams of people with different levels of confidence in their own code.
Programs delete user data, transfer money, make medical decisions, and
deny parole.  Programs are built on the unacknowledged labor of open
source maintainers.  The social and ethical dimensions of programming
are not peripheral concerns to be addressed by external tooling,
governance frameworks, or codes of conduct bolted onto an amoral
substrate.  They are \emph{intrinsic} to the activity of writing
software.

This paper presents \WokeLang{}, a programming language designed
around the premise that ethical deliberation, contribution
attribution, developer confidence, and human readability are
first-class concerns deserving of direct syntactic and semantic
support.  \WokeLang{} does not aim to solve ethics by technical
means---a category error we are keen to
avoid~\citep{green2019good,selbst2019fairness}.  Rather, it aims to
make the social and ethical texture of programming \emph{visible} in
the source code itself, creating affordances for deliberation where
current languages create affordances for speed.

We make the following contributions:

\begin{enumerate}[leftmargin=*]
  \item \textbf{Consent gates} (\S\ref{sec:consent}): a control flow
    construct that creates mandatory deliberation points around
    ethically weighted operations, with formal operational semantics
    that distinguish consent-gated from ungated execution paths.

  \item \textbf{Communicative units of measure}
    (\S\ref{sec:units}): a dimensional type system integrated with the
    consent model, oriented not just toward bug prevention but toward
    making programs legible as descriptions of physical and social
    reality.

  \item \textbf{Gratitude blocks} (\S\ref{sec:gratitude}): an
    executable attribution mechanism that creates auditable graphs of
    contribution acknowledgment within programs.

  \item \textbf{Empathy annotations} (\S\ref{sec:empathy}): a
    metadata system for developer confidence and emotional context,
    surfaceable by code review tools and static analysis.

  \item \textbf{A human-readable design philosophy}
    (\S\ref{sec:philosophy}) that treats syntax as a communication
    medium between people, informed by values-sensitive design and
    feminist HCI.

  \item \textbf{Formal semantics and implementation}
    (\S\ref{sec:semantics}, \S\ref{sec:implementation}): operational
    semantics for the consent gate calculus and a working compiler
    targeting WebAssembly.
\end{enumerate}

% ============================================================================
% 2. BACKGROUND
% ============================================================================
\section{Background and Motivation}
\label{sec:background}

\subsection{Programming Languages as Social Artifacts}

The view that programming languages are purely technical artifacts has
been challenged from multiple directions.  \citet{knuth1984literate}
argued that programs should be written primarily for human readers.
\citet{gabriel1996patterns} explored the relationship between software
design and the built environment through the lens of Christopher
Alexander's pattern language.  More recently, \citet{ko2017past}
surveyed decades of evidence that programming is fundamentally a
\emph{human} activity, shaped by cognitive constraints, social
dynamics, and institutional pressures.

The sociotechnical systems perspective~\citep{trist1981evolution}
insists that technical and social systems are mutually constitutive:
you cannot optimize one without attending to the other.  Applied to
programming languages, this means that syntax and semantics are not
neutral conduits for expressing algorithms; they \emph{afford} certain
ways of thinking and \emph{constrain} others.  A language without
first-class support for error handling affords ignoring errors.  A
language without a concept of ownership affords memory unsafety.  We
contend that a language without a concept of consent affords ethical
inattention.

\subsection{Values-Sensitive Design}

Values-sensitive design (VSD) is a theoretically grounded approach to
the design of technology that accounts for human values in a
principled and comprehensive
manner~\citep{friedman1996value,friedman2013value}.  VSD identifies
three types of investigation: \emph{conceptual} (what values are at
stake?), \emph{empirical} (how do stakeholders experience the
technology?), and \emph{technical} (how can the technology's
properties support or hinder the identified values?).

Applied to programming language design, VSD asks: what values does a
language's syntax make easy to express?  What values does it make
invisible?  Most languages make \emph{efficiency}, \emph{correctness},
and \emph{modularity} easy to express.  Few make \emph{consent},
\emph{attribution}, or \emph{deliberation} expressible at all.

\subsection{Feminist HCI and Care Ethics}

Feminist HCI~\citep{bardzell2010feminist} offers a lens that
foregrounds \emph{pluralism} (whose experience counts?),
\emph{participation} (who gets to shape the design?),
\emph{embodiment} (how does the technology relate to lived
experience?), and \emph{ecology} (how does it fit into broader systems
of power?).  \citet{light2011hci} extends this to consider how design
can ``make space for the expression of values that are currently
marginalized.''

Care ethics~\citep{gilligan1982different,held2006ethics} emphasizes
that moral reasoning is not (only) about abstract principles but about
attending to the concrete needs of others within relationships of
dependency.  Software development is shot through with such
relationships: between developers and users, between maintainers and
contributors, between the people who write code and the people
affected by its execution.

\WokeLang{} draws on these traditions to ask: what would a
programming language look like if it were designed with
\emph{care}---not as an afterthought, but as a structuring principle?

\subsection{Prior Work on Human-Centered Language Design}

Several languages have explored aspects of human-centered design.
Inform~7~\citep{nelson2006natural} uses natural language syntax to
make interactive fiction authoring accessible.  Scratch~\citep{resnick2009scratch}
uses visual blocks to make programming accessible to children.
F\#~\citep{syme2003units} introduced units of measure to the ML
family.  Elm~\citep{czaplicki2012elm} prioritized helpful error
messages.  Pyret~\citep{krishnamurthi2019pyret} integrated testing
into the language itself.

None of these languages, however, have attempted to embed
\emph{ethical deliberation} as a semantic concept.  \WokeLang{} builds
on their insights while pushing into new territory: the idea that
consent, attribution, and emotional context are as worthy of language
support as types, functions, and modules.

% ============================================================================
% 3. CONSENT GATES
% ============================================================================
\section{Consent Gates}
\label{sec:consent}

\subsection{Motivation}

Consider a function that deletes a user's account and all associated
data.  In most programming languages, this function is syntactically
identical to a function that increments a counter.  Both are function
calls; both return values; both can be composed with other functions.
Nothing in the language distinguishes an operation that permanently
destroys someone's personal data from one that adds one to an integer.

This is, we argue, a design failure.  The two operations have
profoundly different \emph{ethical weights}, and a language that treats
them identically provides no affordance for recognizing this
difference.  Developers must rely on comments, naming conventions, or
institutional processes to mark the distinction---all of which are
invisible to the compiler, the type checker, and the runtime.

\WokeLang{} introduces \emph{consent gates} to address this gap.  A
consent gate is a control flow construct that creates a mandatory
deliberation point around an ethically significant operation:

\begin{lstlisting}[caption={Consent gate for account deletion}]
to delete_account(user) ->
  only if okay "Permanently delete all data for {user.name}?" {
    database.drop_user(user.id)
    audit_log.record("deletion", user.id)
    return "Account deleted"
  }
\end{lstlisting}

The \texttt{only if okay} construct is not merely a confirmation
dialog, though it may produce one at runtime.  It is a
\emph{semantic marker} that the enclosed code path involves ethical
weight.  This distinction has consequences at every level of the
system:

\begin{itemize}[leftmargin=*]
  \item \textbf{Compilation:} The compiler tracks consent-gated paths
    separately, enabling analyses such as ``what percentage of
    destructive operations in this codebase are consent-gated?''

  \item \textbf{Type checking:} Functions containing consent gates
    have different types from those without, preventing silent
    composition of gated and ungated operations.

  \item \textbf{Runtime:} The consent mechanism is pluggable---in
    interactive contexts it may prompt a human; in batch contexts it
    may consult a policy file; in test contexts it may auto-approve.

  \item \textbf{Auditing:} Every consent decision (granted or denied)
    is logged with timestamps and the consent string, creating an
    auditable trail of ethical decision points.
\end{itemize}

\subsection{Consent as Effect}

We model consent as an \emph{effect} in the sense of algebraic
effects~\citep{plotkin2009handlers,bauer2015programming}.  A consent
gate introduces the $\consent$ effect, which must be handled by a
consent handler in the calling context.

\begin{definition}[Consent Effect]
\label{def:consent-effect}
The consent effect signature is:
\[
  \consent : \mathsf{String} \to \{\granted, \denied\}
\]
A consent gate \texttt{only if okay $s$ \{ $e$ \}} evaluates $e$
only when the handler for $\consent(s)$ returns $\granted$.
\end{definition}

This formulation has several desirable properties.  First, consent
effects \emph{compose}: a function that calls another consent-gated
function propagates the consent effect upward, ensuring that the
caller is aware of the ethical weight of the callee.  Second, consent
handlers are explicit and auditable: the program must specify how
consent is resolved, whether by user interaction, policy lookup, or
testing mock.  Third, the effect system prevents consent from being
silently swallowed: an unhandled consent effect is a type error.

\subsection{Comparison with Capability-Based Security}

Consent gates share a family resemblance with capability-based
security~\citep{dennis1966programming,miller2006robust}, where access
to a resource is mediated by possession of an unforgeable token.  In
capability systems, the question is ``does this code have the
\emph{right} to perform this operation?''  In consent-gated systems,
the question is ``has a \emph{deliberation} occurred about whether
this operation should happen?''

The distinction matters.  A capability can be granted mechanically: a
policy file says this service may access that database.  Consent
requires \emph{engagement}: someone (a user, a reviewer, a governance
process) must consider the specific operation and its consequences.
Consent gates do not replace capability-based security; they
\emph{complement} it by adding a deliberative layer that capabilities
alone cannot provide.

\begin{table}[t]
\centering
\caption{Comparison of access control paradigms}
\label{tab:access-comparison}
\begin{tabular}{@{}lccc@{}}
\toprule
Property & ACLs & Capabilities & Consent Gates \\
\midrule
Mediates access            & \checkmark & \checkmark & \checkmark \\
Unforgeable tokens         &            & \checkmark &            \\
Requires deliberation      &            &            & \checkmark \\
Auditable decision trail   & Partial    & Partial    & \checkmark \\
Composable (propagates)    &            & \checkmark & \checkmark \\
Type-system integrated     &            & Partial    & \checkmark \\
Human-centered framing     &            &            & \checkmark \\
\bottomrule
\end{tabular}
\end{table}

\subsection{Consent Scoping and Revocation}

Consent in \WokeLang{} is scoped and revocable.  A consent grant
applies only within the lexical scope of the consent gate:

\begin{lstlisting}[caption={Scoped consent}]
only if okay "Process payment of {amount}?" {
  // consent is active here
  charge_card(user.card, amount)
}
// consent has expired; cannot call charge_card
\end{lstlisting}

Consent can also be revoked mid-scope using the \texttt{revoke
consent} construct, modeling situations where circumstances change
during an operation.  This is formalized as a control flow
exception that unwinds the consent-gated block.

\subsection{Nested Consent and Escalation}

When consent gates nest, the inner gate must independently obtain
consent, even if the outer gate has already been approved.  This
models the real-world principle that consenting to one thing does not
imply consenting to everything it entails:

\begin{lstlisting}[caption={Nested consent gates}]
only if okay "Migrate user data to new system?" {
  copy_data(user, new_system)
  only if okay "Delete original data after migration?" {
    delete_original(user, old_system)
  }
}
\end{lstlisting}

The user may consent to migration without consenting to deletion.
This nesting is semantically enforced: the compiler rejects attempts
to hoist inner consent gates out of outer ones.

% ============================================================================
% 4. UNITS OF MEASURE
% ============================================================================
\section{Units of Measure}
\label{sec:units}

\subsection{Motivation: Communication, Not Just Safety}

Units of measure in programming languages are typically motivated by
safety: the Mars Climate Orbiter was lost because of a
Newton-seconds versus pound-force-seconds confusion; the Ariane~5
exploded partly due to a conversion error between 64-bit and 16-bit
representations of a physical quantity~\citep{lions1996ariane}.  These
are compelling reasons to build dimensional analysis into type systems.

\WokeLang{} shares this safety motivation but extends it with a
\emph{communicative} one.  Units of measure are not just constraints
that prevent bugs; they are \emph{documentation} that makes programs
legible as descriptions of physical and social reality.  When a
programmer writes:

\begin{lstlisting}[caption={Units as communication}]
remember distance = 42 measured in km
remember time = 3.5 measured in hours
remember speed = distance / time  // inferred: km/hours
\end{lstlisting}

the units serve as a form of \emph{self-narration}: the code explains
not just what it computes but what the computation \emph{means} in
physical terms.  This aligns with \citeauthor{knuth1984literate}'s
vision of literate programming, extended from prose annotations to the
type system itself.

\subsection{The Dimensional Type System}

\WokeLang{}'s type system extends Hindley-Milner inference with
dimensional analysis following \citet{kennedy1997relational}.  Each
numeric type carries an associated \emph{dimension}, which is an
element of the free abelian group generated by base dimensions.

\begin{definition}[Dimensions]
\label{def:dimensions}
Let $\mathcal{B} = \{L, M, T, I, \Theta, N, J, \ldots\}$ be a set of
base dimensions (length, mass, time, current, temperature, amount,
luminous intensity, and user-defined dimensions).  A \emph{dimension}
is an element of the free abelian group $\mathbb{Z}^{|\mathcal{B}|}$,
written multiplicatively:
\[
  d = L^{a_1} \cdot M^{a_2} \cdot T^{a_3} \cdots
\]
The dimension of a dimensionless quantity is the identity element
$\mathbf{1}$.
\end{definition}

\begin{definition}[Dimensioned Types]
\label{def:dim-types}
A \emph{dimensioned type} is a pair $\unitty{\tau, d}$ where $\tau$
is a numeric type (\texttt{Int}, \texttt{Float}, \texttt{Decimal})
and $d$ is a dimension.  The typing rules for arithmetic operations
respect dimensional algebra:
\begin{mathpar}
  \inferrule{
    \Gamma \vdash e_1 : \unitty{\tau, d_1} \\
    \Gamma \vdash e_2 : \unitty{\tau, d_2}
  }{
    \Gamma \vdash e_1 + e_2 : \unitty{\tau, d_1}
  }\;\text{if } d_1 = d_2
  \and
  \inferrule{
    \Gamma \vdash e_1 : \unitty{\tau, d_1} \\
    \Gamma \vdash e_2 : \unitty{\tau, d_2}
  }{
    \Gamma \vdash e_1 \times e_2 : \unitty{\tau, d_1 \cdot d_2}
  }
  \and
  \inferrule{
    \Gamma \vdash e_1 : \unitty{\tau, d_1} \\
    \Gamma \vdash e_2 : \unitty{\tau, d_2}
  }{
    \Gamma \vdash e_1 \div e_2 : \unitty{\tau, d_1 \cdot d_2^{-1}}
  }
\end{mathpar}
\end{definition}

\subsection{User-Defined Dimensions}

Beyond SI base units, \WokeLang{} allows user-defined dimensions for
domain-specific quantities:

\begin{lstlisting}[caption={User-defined dimensions}]
type Currency is dimension
type UserCount is dimension

remember revenue = 50000 measured in Currency
remember users = 1200 measured in UserCount
remember arpu = revenue / users  // Currency/UserCount
\end{lstlisting}

This is particularly useful in domains where informal ``units'' are
pervasive but rarely tracked: dollars per user, requests per second,
students per classroom.  Making these units explicit turns tribal
knowledge into type-checked documentation.

\subsection{Integration with Consent Gates}

The dimensional type system interacts with consent gates in a subtle
way.  Certain unit conversions may themselves be consent-gated,
particularly when the conversion involves a loss of precision or a
change of meaning:

\begin{lstlisting}[caption={Consent-gated conversion}]
remember precise_value = 3.14159265 measured in radians
// Lossy conversion requires consent
only if okay "Convert to degrees (loses precision)?" {
  remember approx = precise_value as degrees
}
\end{lstlisting}

This models the real-world principle that unit conversions are not
always innocent: converting between coordinate systems, currency
denominations, or medical units can introduce errors that have ethical
consequences.

% ============================================================================
% 5. GRATITUDE BLOCKS
% ============================================================================
\section{Gratitude Blocks}
\label{sec:gratitude}

\subsection{The Attribution Problem}

Open source software depends on the unpaid labor of maintainers whose
contributions are systematically
undervalued~\citep{eghbal2020working}.  Attribution in current
practice is handled through comments, \texttt{CONTRIBUTORS} files,
commit histories, and social conventions---none of which are visible
to the compiler, queryable at runtime, or auditable by automated
tools.

\WokeLang{} introduces \emph{gratitude blocks} to make attribution a
first-class, executable construct:

\begin{lstlisting}[caption={Gratitude block}]
thanks to {
  "Alice Chen"       -> "Designed the consent protocol"
  "Bob Martinez"     -> "Fixed the race condition in gate resolution"
  "Fatima Al-Rashid" -> "Reviewed the formal semantics"
}

to resolve_consent(gate) ->
  // implementation that benefited from all three contributions
  ...
\end{lstlisting}

\subsection{Semantics of Gratitude}

A gratitude block is not a comment.  It is a program construct with
runtime semantics: when a gratitude block is evaluated, it emits
\emph{attribution events} that are recorded in the program's
attribution graph.

\begin{definition}[Attribution Graph]
\label{def:attr-graph}
The attribution graph $G = (V, E)$ is a directed graph where:
\begin{itemize}
  \item Vertices $V = V_P \cup V_C$ are partitioned into
    \emph{person vertices} $V_P$ (contributors) and \emph{code
    vertices} $V_C$ (modules, functions, blocks).
  \item Edges $E \subseteq V_P \times V_C \times \Sigma^*$ are
    labeled with attribution descriptions.
\end{itemize}
The graph is accumulated during program execution and can be queried
at any point via the standard library:
\end{definition}

\begin{lstlisting}[caption={Querying the attribution graph}]
remember contributors = attribution.query("resolve_consent")
for each person in contributors {
  say("{person.name}: {person.contribution}")
}
\end{lstlisting}

\subsection{Attribution Propagation}

When a function $f$ calls a function $g$ that contains a gratitude
block, the attribution edges from $g$ are transitively visible from
$f$.  This models the reality that when you use a library, you
benefit from the work of its contributors even if you never see their
names:

\begin{lstlisting}[caption={Transitive attribution}]
// In module authentication.woke
thanks to { "Priya Sharma" -> "Implemented TOTP" }
to verify_totp(code, secret) -> ...

// In module login.woke
to login(user, code) ->
  // Priya's contribution is transitively attributed
  if verify_totp(code, user.secret) { ... }

// Query: attribution.transitive("login") includes Priya
\end{lstlisting}

\subsection{Social and Institutional Implications}

The attribution graph has implications beyond individual programs.
Organizations can aggregate attribution data across projects to
understand contribution patterns, identify unacknowledged labor, and
make invisible work visible.  Package managers can surface attribution
information alongside dependency trees, so that \texttt{install
http-server} also shows you who built the foundations you depend on.

This is, explicitly, a \emph{political} design choice.  We believe
that programming languages should make it easy to acknowledge the work
of others, not because acknowledgment is technically necessary (the
code runs either way) but because it is \emph{ethically necessary} in
a field that systematically undervalues maintenance, documentation,
and invisible labor~\citep{star1999ethnography,eghbal2020working}.

% ============================================================================
% 6. EMPATHY ANNOTATIONS
% ============================================================================
\section{Empathy Annotations}
\label{sec:empathy}

\subsection{The Confidence Gap}

Code review is a social process~\citep{rigby2013convergent}.
Reviewers attempt to understand not just \emph{what} the code does but
\emph{how confident the author is} in their implementation.  This
confidence information is typically conveyed through social signals:
the tone of the pull request description, the number of ``TODO''
comments, the developer's reputation, informal conversations.  None of
this is captured in the code itself.

\WokeLang{} introduces \emph{empathy annotations}, a metadata system
that allows developers to attach confidence and emotional context
directly to code:

\begin{lstlisting}[caption={Empathy annotations}]
@feeling(confident=true)
to add(a, b) -> a + b

@feeling(confident=false, note="Edge case handling is fragile")
to parse_date(input) ->
  // I think this handles leap years correctly...
  ...

@feeling(anxious=true, note="First time writing concurrent code")
to worker_pool(size) ->
  ...
\end{lstlisting}

\subsection{Annotation Semantics}

Empathy annotations are structured metadata, not free-form comments.
They have a defined schema:

\begin{lstlisting}[caption={Empathy annotation schema}]
type Feeling is {
  confident: Boolean       // author's confidence in correctness
  anxious: Boolean         // author's anxiety about the code
  experimental: Boolean    // this is exploratory/prototype code
  needs_review: Boolean    // explicitly requesting review
  note: String             // free-form context
}
\end{lstlisting}

Because the schema is fixed, empathy annotations are \emph{queryable}
by tooling.  A code review tool can:

\begin{itemize}[leftmargin=*]
  \item Surface all functions marked \texttt{confident=false} for
    priority review.
  \item Highlight code marked \texttt{anxious=true} with additional
    context about the author's concerns.
  \item Track the ratio of confident to unconfident code as a project
    health metric.
  \item Generate reports showing where uncertainty concentrates in the
    codebase.
\end{itemize}

\subsection{Making the Implicit Explicit}

The design of empathy annotations is motivated by a simple
observation: developers already have feelings about their code.  They
already know which functions are solid and which are held together with
hope and string.  But this knowledge is locked in their heads, shared
only through informal channels that disappear when the developer
leaves the project.

By making this knowledge explicit and structured, empathy annotations
create what \citet{star1999ethnography} calls
\emph{infrastructure}---the boring, invisible substrate that makes
coordination possible.  The annotation \texttt{@feeling(confident=%
false)} is a tiny act of vulnerability that enables a more honest,
more efficient review process.

\subsection{Potential for Misuse}

We acknowledge the potential for empathy annotations to be misused.
An employer might penalize developers who mark code as unconfident.  A
culture of toxic positivity might pressure developers to mark
everything as confident.  These are real risks, and they are not
solvable by technical means alone.

Our position is that the \emph{option} to express uncertainty is
better than the absence of that option.  A language that provides no
way to say ``I am not sure about this'' is not neutral; it implicitly
demands confidence, or at least the performance of confidence.
\WokeLang{} provides the affordance; institutional and cultural
practices must ensure it is used well.  We discuss this further in
\S\ref{sec:discussion}.

% ============================================================================
% 7. THE HUMAN-READABLE DESIGN PHILOSOPHY
% ============================================================================
\section{The Human-Readable Design Philosophy}
\label{sec:philosophy}

\subsection{Language as Communication Medium}

\WokeLang{}'s syntax is designed to read like prose.  This is not a
superficial aesthetic choice; it reflects a specific theory of what
programming languages \emph{are}.

Most programming language design treats syntax as a notation for
expressing computations to a machine, optimized for precision and
conciseness.  This is a reasonable design goal, but it is not the only
one.  \WokeLang{} treats syntax as a \emph{communication medium}
between humans, optimized for readability, clarity, and the expression
of intent.

\begin{lstlisting}[caption={Human-readable syntax examples}]
// Function definition reads as English prose
to calculate_shipping(weight, destination) ->
  remember base_cost = 5.99 measured in USD
  remember per_kg = 2.50 measured in USD/kg
  return base_cost + (weight * per_kg)

// Pipelines read left to right
remember result = raw_data
  then validate
  then normalize
  then analyze

// Pattern matching reads as case analysis
match user.role with
  when "admin"     -> grant_full_access(user)
  when "moderator" -> grant_mod_access(user)
  otherwise        -> grant_read_access(user)
\end{lstlisting}

\subsection{Design Principles}

The human-readable design philosophy is guided by several principles:

\paragraph{Principle 1: Keywords as Verbs.}
Where most languages use symbols (\texttt{=>}, \texttt{|>},
\texttt{::}), \WokeLang{} prefers English words (\texttt{then},
\texttt{when}, \texttt{measured in}).  This slightly increases
verbosity in exchange for dramatically increased readability for
non-experts and for future maintainers encountering the code for the
first time.

\paragraph{Principle 2: Declarations as Statements of Intent.}
\texttt{remember x = 5} is more evocative than \texttt{let x = 5} or
\texttt{var x = 5}.  The word ``remember'' frames variable binding
not as a machine operation (allocating a slot) but as a cognitive act
(committing a fact to memory).  The word \texttt{to} before function
definitions frames the function as a \emph{purpose}: ``to add $a$ and
$b$'' rather than ``function add.''

\paragraph{Principle 3: Pipelines as Narratives.}
The \texttt{then} keyword creates data pipelines that read as
narratives: ``take the raw data, \emph{then} validate it, \emph{then}
normalize it, \emph{then} analyze it.''  This mirrors how humans
describe processes in everyday language, reducing the cognitive
distance between a natural-language description and its implementation.

\paragraph{Principle 4: Consent as Conversation.}
The \texttt{only if okay} construct frames potentially destructive
operations as conversations between the program and its operator.
This is deliberate: the phrasing invites the reader to imagine the
ethical deliberation that should accompany the operation.

\subsection{Readability as Ethical Choice}

We argue that readability is not merely a convenience but an
\emph{ethical} property of programming languages.  Code that is
difficult to read is code that is difficult to review, difficult to
audit, and difficult to hold accountable.  Obscure code concentrates
power in the hands of those who can decode it.  Readable code
distributes power more broadly, enabling more people to participate in
the scrutiny and governance of software systems.

This position is informed by \citet{nissenbaum2004privacy}'s concept
of \emph{contextual integrity}: the ethical assessment of a technology
depends not just on what it does but on how it fits into existing
social contexts and norms.  A programming language that is legible
only to specialists creates a context in which software governance is
delegated to specialists.  A language that is legible to a broader
audience creates a context in which governance can be more broadly
shared.

% ============================================================================
% 8. FORMAL SEMANTICS
% ============================================================================
\section{Formal Semantics}
\label{sec:semantics}

We now provide operational semantics for the consent gate calculus,
which is the most technically novel aspect of \WokeLang{}'s design.

\subsection{Syntax}

We define a core calculus $\lambda_{\consent}$ that extends the simply
typed lambda calculus with consent gates:

\begin{definition}[Syntax of $\lambda_{\consent}$]
\label{def:syntax}
\begin{align*}
  e ::=\; & x \mid \lambda x:\tau.\, e \mid e_1\; e_2 \mid n \mid
    e_1 + e_2 \mid \\
    & \mathbf{only\_if\_okay}\; s\; e \mid
    \mathbf{consent\_handler}\; h\; e \mid \\
    & \mathbf{attr}\; p\; d\; e \mid
    \mathbf{feel}\; f\; e \\[0.5em]
  v ::=\; & \lambda x:\tau.\, e \mid n \mid
    \mathbf{granted} \mid \mathbf{denied} \\[0.5em]
  \tau ::=\; & \mathsf{Int} \mid \mathsf{String} \mid
    \mathsf{Bool} \mid \tau_1 \to \tau_2 \mid
    \unitty{\tau, d} \mid \\
    & \mathsf{Consent}[\tau] \\[0.5em]
  h ::=\; & \mathsf{interactive} \mid
    \mathsf{policy}(P) \mid
    \mathsf{auto}(\granted) \mid
    \mathsf{auto}(\denied)
\end{align*}
\end{definition}

The key new forms are $\mathbf{only\_if\_okay}\; s\; e$, which
creates a consent gate with prompt string $s$ guarding expression $e$;
and $\mathbf{consent\_handler}\; h\; e$, which installs a consent
handler $h$ for the duration of $e$.

\subsection{Typing Rules}

\begin{definition}[Typing for consent gates]
\label{def:consent-typing}
We extend the typing judgment $\Gamma; \Sigma \vdash e : \tau$ with a
consent context $\Sigma$ that tracks whether we are inside a consent
handler.

\begin{mathpar}
  \inferrule[\textsc{T-Consent}]{
    \Gamma; \Sigma \vdash e : \tau \\
    s : \mathsf{String}
  }{
    \Gamma; \Sigma \vdash
    \mathbf{only\_if\_okay}\; s\; e : \mathsf{Consent}[\tau]
  }
  \and
  \inferrule[\textsc{T-Handler}]{
    \Gamma; \Sigma, h \vdash e : \tau
  }{
    \Gamma; \Sigma \vdash
    \mathbf{consent\_handler}\; h\; e : \tau
  }
  \and
  \inferrule[\textsc{T-Attr}]{
    \Gamma; \Sigma \vdash e : \tau \\
    p : \mathsf{String} \\
    d : \mathsf{String}
  }{
    \Gamma; \Sigma \vdash
    \mathbf{attr}\; p\; d\; e : \tau
  }
  \and
  \inferrule[\textsc{T-Feel}]{
    \Gamma; \Sigma \vdash e : \tau \\
    f : \mathsf{Feeling}
  }{
    \Gamma; \Sigma \vdash
    \mathbf{feel}\; f\; e : \tau
  }
\end{mathpar}
\end{definition}

The type $\mathsf{Consent}[\tau]$ wraps the result type to indicate
that the value may not be available (if consent is denied).  This
forces callers to handle the possibility of denial.

\subsection{Operational Semantics}

\begin{definition}[Small-step semantics for consent gates]
\label{def:consent-opsem}
We define the reduction relation $\langle e, H, A \rangle \step
\langle e', H', A' \rangle$ where $H$ is the consent handler stack
and $A$ is the attribution accumulator.

\begin{mathpar}
  \inferrule[\textsc{E-ConsentGranted}]{
    H = h :: H' \\
    h(s) = \granted
  }{
    \langle \mathbf{only\_if\_okay}\; s\; e,\; H,\; A \rangle
    \step
    \langle e,\; H,\; A \cdot \mathsf{log}(\granted, s, t) \rangle
  }
  \and
  \inferrule[\textsc{E-ConsentDenied}]{
    H = h :: H' \\
    h(s) = \denied
  }{
    \langle \mathbf{only\_if\_okay}\; s\; e,\; H,\; A \rangle
    \step
    \langle \denied,\; H,\; A \cdot \mathsf{log}(\denied, s, t) \rangle
  }
  \and
  \inferrule[\textsc{E-Handler}]{
    \langle e,\; h :: H,\; A \rangle \step
    \langle e',\; h :: H,\; A' \rangle
  }{
    \langle \mathbf{consent\_handler}\; h\; e,\; H,\; A \rangle
    \step
    \langle \mathbf{consent\_handler}\; h\; e',\; H,\; A' \rangle
  }
  \and
  \inferrule[\textsc{E-HandlerDone}]{
    v \text{ is a value}
  }{
    \langle \mathbf{consent\_handler}\; h\; v,\; H,\; A \rangle
    \step
    \langle v,\; H,\; A \rangle
  }
  \and
  \inferrule[\textsc{E-Attr}]{
    \langle e,\; H,\; A \cdot (p, d) \rangle \step
    \langle e',\; H,\; A' \rangle
  }{
    \langle \mathbf{attr}\; p\; d\; e,\; H,\; A \rangle \step
    \langle e',\; H,\; A' \rangle
  }
  \and
  \inferrule[\textsc{E-Feel}]{
    \langle e,\; H,\; A \rangle \step
    \langle e',\; H,\; A' \rangle
  }{
    \langle \mathbf{feel}\; f\; e,\; H,\; A \rangle \step
    \langle e',\; H,\; A' \rangle
  }
\end{mathpar}
\end{definition}

Note that the \textsc{E-Feel} rule is an identity on the computation:
empathy annotations do not affect the operational behavior of the
program.  They are ``phantom'' metadata---present in the source,
erased at runtime, but available to static analysis tools.  The
attribution rule \textsc{E-Attr} accumulates attribution edges in $A$
but similarly does not affect the computational result.

\subsection{Properties}

\begin{theorem}[Progress]
\label{thm:progress}
If $\Gamma; \Sigma \vdash e : \tau$ and $e$ is not a value, then
there exists $e'$, $H'$, $A'$ such that $\langle e, H, A \rangle
\step \langle e', H', A' \rangle$, provided $H$ contains a handler
for any consent effect in $e$.
\end{theorem}

\begin{proof}
By structural induction on the typing derivation.  The critical case
is \textsc{T-Consent}: we strengthen the paper rule to require
$\Sigma \neq \varnothing$ (at least one handler in scope), closing a
gap in the original sketch.  With this invariant, the top handler $h$
is always available and total, so one of \textsc{E-ConsentGranted} or
\textsc{E-ConsentDenied} fires.  The full machine-checked proof is in
\texttt{src/abi/WokeLang/Safety.idr} (\texttt{progress});
see the supplementary proof inventory
(\texttt{docs/supplementary/safety-proofs.adoc}) for details.
\end{proof}

\begin{theorem}[Preservation]
\label{thm:preservation}
If $\Gamma; \Sigma \vdash e : \tau$ and $\langle e, H, A \rangle
\step \langle e', H', A' \rangle$, then there exists $\tau'$ with
$\Gamma; \Sigma' \vdash e' : \tau'$, where $\tau' = \tau$ for most
rules, $\tau' = \tau$ (body type) when \textsc{E-ConsentGranted}
fires, and $\tau' = \mathsf{TyDenied}$ when \textsc{E-ConsentDenied}
fires.
\end{theorem}

\begin{proof}
By structural induction on the reduction derivation.  The consent
cases follow from handler totality: \textsc{E-ConsentGranted} returns
$e$ (the body, typed at $\tau$), and \textsc{E-ConsentDenied} returns
\texttt{EDenied} (typed at $\mathsf{TyDenied}$).  The E-Beta case
uses a closed-value substitution lemma proved by induction with a
context-exchange step for nested lambda binders.  Machine-checked in
\texttt{src/abi/WokeLang/Safety.idr} (\texttt{preservation}).
\end{proof}

\begin{theorem}[Consent Non-Evasion]
\label{thm:non-evasion}
In a well-typed $\lambda_{\consent}$ program, it is not possible to
evaluate the body of a consent gate without the corresponding consent
handler producing either $\granted$ or $\denied$ and logging the
decision.
\end{theorem}

\begin{proof}
By inversion on the \textsc{Step} derivation for $\mathbf{only\_if\_okay}\;s\;e$.
The only constructors of \textsc{Step} whose head expression matches
$\mathbf{only\_if\_okay}$ are \textsc{E-ConsentGranted} and
\textsc{E-ConsentDenied}; all other constructors match different head
forms and are structurally absent.  Both firing rules require
$H = h :: H'$ (non-empty stack) and append
$\mathsf{ConsentEntry}(d, s)$ to $A$.
Machine-checked in \texttt{src/abi/WokeLang/Safety.idr}
(\texttt{nonEvasion}, \texttt{consentBodyRequiresLog}).
\end{proof}

This theorem is the core safety property of the consent system: you
cannot execute ethically weighted code without going through the
deliberation step.

% ============================================================================
% 9. IMPLEMENTATION
% ============================================================================
\section{Implementation}
\label{sec:implementation}

\subsection{Architecture}

\WokeLang{} is implemented as a multi-stage compiler with the
following pipeline:

\begin{enumerate}
  \item \textbf{Lexer} (Rust, using \texttt{logos}): Tokenizes
    \WokeLang{} source files into a stream of typed tokens, handling
    the natural-language keywords (\texttt{remember}, \texttt{to},
    \texttt{only if okay}, etc.) and empathy annotations.

  \item \textbf{Parser} (Rust, hand-written recursive descent):
    Produces an abstract syntax tree (AST) that includes consent gate
    nodes, gratitude blocks, unit annotations, and empathy metadata as
    first-class AST nodes (not comments or pragmas).

  \item \textbf{Type Checker} (Rust): Implements Hindley-Milner type
    inference extended with dimensional analysis and consent effects.
    The type checker maintains a consent context stack and rejects
    programs with unhandled consent effects.

  \item \textbf{Bytecode Compiler} (Rust): Compiles the typed AST
    to a custom bytecode format (\texttt{.wbc}) with dedicated
    opcodes for consent gates, attribution emission, and unit
    propagation.

  \item \textbf{Virtual Machine} (Rust): A stack-based VM that
    executes the bytecode, with a pluggable consent handler interface
    and a built-in attribution accumulator.

  \item \textbf{WebAssembly Backend} (planned): Compilation to WASM
    for browser-based execution, with consent handlers bridged to the
    host environment via the WASM component model.
\end{enumerate}

\subsection{Consent Handler Interface}

The runtime provides a trait-based consent handler interface:

\begin{lstlisting}[language=Rust,caption={Consent handler trait (Rust)}]
/// Consent handler determines how consent gates are resolved.
pub trait ConsentHandler {
    /// Called when a consent gate is encountered.
    /// Returns true (granted) or false (denied).
    fn request_consent(&mut self, prompt: &str) -> bool;

    /// Called after consent resolution for audit logging.
    fn log_decision(
        &self, prompt: &str, granted: bool, timestamp: u64
    );
}

pub struct InteractiveHandler;  // prompts user via stdin
pub struct PolicyHandler { policy: PolicyFile; }  // consults file
pub struct TestHandler { responses: Vec<bool>; }  // for testing
\end{lstlisting}

This design allows the same \WokeLang{} program to run in different
consent contexts without modification: interactively during
development, policy-driven in production, auto-approving in tests.

\subsection{Attribution Accumulator}

The attribution accumulator is implemented as an append-only log that
records every attribution event during program execution.  It can be
serialized to JSON, queried at runtime through the standard library,
and aggregated across multiple program runs.

\begin{lstlisting}[language=Rust,caption={Attribution event structure}]
/// A single attribution event in the program trace.
pub struct AttributionEvent {
    pub contributor: String,
    pub description: String,
    pub source_location: SourceSpan,
    pub timestamp: u64,
    pub transitive_from: Option<Vec<AttributionEvent>>,
}
\end{lstlisting}

\subsection{Empathy Annotations in the Toolchain}

Empathy annotations are preserved in the AST but erased during
bytecode compilation (they have no runtime semantics).  They are
available to:

\begin{itemize}[leftmargin=*]
  \item The \textbf{LSP server}, which can display confidence
    information inline in the editor.
  \item The \textbf{code review integration}, which surfaces
    \texttt{confident=false} annotations in pull request views.
  \item The \textbf{project health dashboard}, which aggregates
    empathy metadata across the codebase to show confidence
    distribution.
  \item The \textbf{documentation generator}, which can include
    empathy information in API docs (``the author notes that edge case
    handling in this function is fragile'').
\end{itemize}

\subsection{Performance Considerations}

Consent gates introduce a potential runtime cost: each gate invokes
the consent handler, which may involve I/O (user prompts, policy file
reads, network calls).  We mitigate this through:

\begin{itemize}[leftmargin=*]
  \item \textbf{Consent caching:} Within a single execution, repeated
    consent requests with the same prompt can be cached (configurable
    per handler).
  \item \textbf{Batch consent:} Multiple consent gates can be
    collected and presented as a single batch prompt.
  \item \textbf{Compile-time consent elimination:} When the consent
    handler is statically known (e.g., \texttt{auto(granted)} in test
    mode), the compiler can eliminate the consent gate entirely.
\end{itemize}

Attribution accumulation has negligible overhead (append-only log).
Empathy annotations have zero runtime cost (erased at compilation).
Unit checking is entirely compile-time; the runtime representation of
\texttt{5 measured in km} is simply the integer \texttt{5}.

% ============================================================================
% 10. DISCUSSION
% ============================================================================
\section{Discussion}
\label{sec:discussion}

\subsection{Should Programming Languages Encode Values?}

The most fundamental objection to \WokeLang{} is that programming
languages should be \emph{value-neutral}---that they should provide
general-purpose computation facilities and leave ethical questions to
the humans using them.

We reject this framing on both descriptive and normative grounds.
\emph{Descriptively}, no programming language is value-neutral.  C
values programmer autonomy over safety.  Rust values memory safety
over ease of use.  Python values readability over performance.  These
are not accidental properties; they are conscious design
choices that shape how millions of people write software.  The
question is not whether a language embeds values, but \emph{which}
values it embeds and whether it does so \emph{transparently}.

\emph{Normatively}, we argue that language designers have a
responsibility to consider the ethical implications of their design
choices, just as architects have a responsibility to consider the
social implications of building
design~\citep{winner1980artifacts,flanagan2005values}.  A language
that makes it syntactically trivial to delete user data and
syntactically impossible to acknowledge the contributor who wrote the
deletion safeguards has made an ethical choice---it has just made it
invisibly.

\subsection{The Risk of Ethical Theater}

A legitimate concern is that consent gates, gratitude blocks, and
empathy annotations could become empty rituals---box-checking
exercises that provide the \emph{appearance} of ethical deliberation
without the substance.  A developer might write \texttt{only if okay
"Delete?"} and then configure the consent handler to auto-approve
everything.

This risk is real, but it is not unique to \WokeLang{}.  Type systems
can be circumvented with unsafe casts.  Tests can be written to pass
without testing anything.  Code reviews can be rubber-stamped.  The
existence of a mechanism does not guarantee its good-faith use; it
only \emph{creates the possibility}.

What \WokeLang{} provides that comments and institutional processes do
not is \emph{auditability}.  Because consent decisions are logged,
auto-approve configurations are \emph{visible} in the audit trail.
Because gratitude blocks are executable, their absence is detectable
by tooling.  Because empathy annotations have a fixed schema, their
distribution across a codebase is measurable.  These properties do not
prevent bad faith, but they make it harder to sustain invisibly.

\subsection{Cultural and Linguistic Considerations}

\WokeLang{}'s English keywords raise questions about linguistic
inclusion.  The choice to use English words like ``remember,'' ``only
if okay,'' and ``thanks to'' privileges English speakers.  We
acknowledge this limitation and note that it is shared by virtually
every mainstream programming language.

Future work could explore localization of keywords (as
AppleScript once attempted) or a keyword-agnostic intermediate
representation that could be presented in different natural languages.
We consider this important but orthogonal to the core semantic
contributions of this paper.

\subsection{Relationship to Regulation}

\WokeLang{}'s consent gates have a natural affinity with regulatory
frameworks like the GDPR's requirement for informed
consent~\citep{gdpr2016} and the EU AI Act's requirements for
human oversight of high-risk AI
systems~\citep{euaiact2024}.  A consent-gated
codebase provides built-in evidence of deliberation that may be
relevant to regulatory compliance.

We caution, however, against treating consent gates as a
\emph{compliance tool}.  Their purpose is to create affordances for
ethical deliberation, not to generate audit trails for regulators.
If consent gates are adopted primarily for compliance, they risk
becoming exactly the kind of empty ritual we warned against in the
previous section.

\subsection{Scalability of Consent}

A practical concern is \emph{consent fatigue}: if too many operations
are consent-gated, developers will habituate to the prompts and
approve without deliberation.  This mirrors the well-known
``alert fatigue'' problem in
security~\citep{anderson2008security,akhawe2013alice}.

\WokeLang{} mitigates this through design guidance: consent gates
should be reserved for operations with genuine ethical weight, not for
routine computations.  The type system helps enforce this by making
consent a visible cost in function signatures, creating social
pressure to use it judiciously.  We also note that consent fatigue is
a problem with \emph{overuse}, not with the concept itself; the
appropriate number of consent gates in a program is a design judgment,
not a technical parameter.

% ============================================================================
% 11. RELATED WORK
% ============================================================================
\section{Related Work}
\label{sec:related}

\subsection{Units of Measure in Type Systems}

F\#'s units of measure~\citep{syme2003units,kennedy1997relational}
are the most mature implementation of dimensional analysis in a
mainstream language.  \WokeLang{} follows F\#'s approach in modeling
dimensions as elements of a free abelian group but differs in
treating units as primarily \emph{communicative} rather than
primarily \emph{preventive}.  Frink~\citep{frink2002} takes an
even more aggressive approach, making unit conversion a fundamental
operation.

\subsection{Natural Language Programming}

Inform~7~\citep{nelson2006natural} demonstrated that natural language
syntax can work in a production language, at least for the domain of
interactive fiction.  AppleScript~\citep{cook2007applescript}
attempted English-like syntax for scripting but was widely criticized
for being neither natural enough to read easily nor formal enough to
write without consulting documentation.  \WokeLang{}'s approach is
more modest: it uses English keywords where they aid readability
(\texttt{remember}, \texttt{then}, \texttt{only if okay}) but retains
conventional syntax for expressions and types.

\subsection{Values-Sensitive Design in Software}

\citet{friedman2013value} provide the theoretical framework for VSD.
\citet{flanagan2005values} specifically address values in
design and argue that technology is never value-neutral.
\citet{nissenbaum2004privacy} introduces contextual integrity as a
framework for evaluating privacy-affecting technologies.
\citet{winner1980artifacts} famously argued that artifacts have
politics.  \WokeLang{} is, to our knowledge, the first programming
language to explicitly engage with this literature in its design.

\subsection{Feminist HCI}

\citet{bardzell2010feminist} established the field of feminist HCI.
\citet{light2011hci} extended it to consider how design can challenge
existing power structures.  \citet{dantec2009infrastructuring}
explored how design can create the conditions for participation.
\WokeLang{}'s emphasis on readability, attribution, and the
expression of uncertainty draws on these perspectives.

\subsection{Ethics in Computing}

\citet{green2019good} warns against ``ethics-washing'' in AI, where
ethical language is adopted without substantive change.
\citet{selbst2019fairness} identifies ``abstraction traps'' in
fairness research, where formal definitions of fairness fail to
capture social context.  \WokeLang{} takes these warnings seriously:
consent gates are not a solution to ethical problems in computing,
but an \emph{affordance} that makes ethical deliberation cheaper and
more visible.

\subsection{Capability-Based Security}

The capability model~\citep{dennis1966programming,miller2006robust}
mediates access through unforgeable tokens.  Object-capability
languages like E~\citep{miller2006robust} and
Pony~\citep{clebsch2015deny} integrate capabilities into the type
system.  Consent gates complement capabilities by adding a
deliberative dimension that pure capability systems lack (see
\S\ref{sec:consent}).

\subsection{Algebraic Effects and Handlers}

The consent effect system draws on algebraic effects and
handlers~\citep{plotkin2009handlers,bauer2015programming}, as
implemented in languages like
Eff~\citep{bauer2012programming},
Koka~\citep{leijen2017type}, and
Multicore OCaml~\citep{sivaramakrishnan2021retrofitting}.  Our
contribution is not a new effect system but a novel \emph{application}
of effects to model ethical deliberation.

\subsection{Literate Programming and Documentation}

\citet{knuth1984literate} introduced literate programming as the
idea that programs should be written for human readers.
\citet{gabriel1996patterns} explored the relationship between
software and architecture.  Doctest systems in
Python~\citep{doctest2004} and Rust blur the line between
documentation and executable code.  \WokeLang{}'s gratitude blocks
and empathy annotations push further: they make social metadata
(attribution, confidence) part of the program itself, not just its
documentation.

\subsection{Emotional and Social Aspects of Programming}

\citet{graziotin2014happy} found that happy developers are more
productive.  \citet{muller2015stuck} studied how developers experience
and cope with being stuck.  \citet{ford2015paradise} examined how
affective states influence debugging.  Empathy annotations make these
well-studied emotional dimensions of programming \emph{visible} in
the source code, creating a bridge between the social science of
software development and the artifacts it studies.

% ============================================================================
% 12. FUTURE WORK
% ============================================================================
\section{Future Work}
\label{sec:future}

Several directions warrant further investigation:

\paragraph{Consent Delegation Protocols.}
In large systems, consent decisions may need to be delegated across
organizational boundaries.  We plan to explore formal protocols for
consent delegation, drawing on work in distributed authorization.

\paragraph{Attribution and Supply Chain Security.}
The attribution graph could be extended to support software supply
chain integrity, connecting gratitude blocks to SBOM (Software Bill of
Materials) generation and provenance tracking.

\paragraph{Empirical Evaluation.}
We have not yet conducted user studies to evaluate whether consent
gates, gratitude blocks, and empathy annotations improve ethical
deliberation, attribution practices, or code review quality in
practice.  Controlled experiments with professional developers are a
priority.

\paragraph{Formal Verification.}
The $\lambda_{\consent}$ calculus presented here covers a simplified
core language.  Proofs of Progress, Preservation, and Non-Evasion have
been mechanised in Idris~2~\citep{brady2021idris2} under
\texttt{\%default total} with no uses of \texttt{believe\_me} or
\texttt{postulate}; see \texttt{src/abi/WokeLang/Safety.idr} and
\texttt{docs/supplementary/safety-proofs.adoc}.  The mechanisation
strengthens \textsc{T-Consent} to require a non-empty handler context
$\Sigma$, fixing a gap in the original proof sketch.  Extending the
formal treatment to the full language---including higher-order
functions, mutation, concurrency, and module systems---remains future
work.

\paragraph{Localization.}
Exploring multilingual keyword support while maintaining semantic
precision is an open design challenge.

\paragraph{Integration with Existing Ecosystems.}
\WokeLang{}'s consent gates and attribution system could be exposed
as a library or language extension for existing languages (e.g., a
Rust procedural macro, an OCaml PPX extension), lowering the barrier
to adoption.

% ============================================================================
% 13. CONCLUSION
% ============================================================================
\section{Conclusion}
\label{sec:conclusion}

Programming languages are not neutral.  They embody assumptions about
what matters, and those assumptions shape the software that billions
of people depend on.  For fifty years, the dominant assumption has
been that what matters is computation---correctness, performance,
modularity.  These are real and important values, and we do not
propose to abandon them.

But they are not the only values that matter.  Consent matters:
software should not perform ethically significant actions without
deliberation.  Attribution matters: the people who build and maintain
software deserve acknowledgment.  Confidence matters: the implicit
social signals that pervade code review should be made explicit and
queryable.  Readability matters: code is a communication medium
between humans, and clarity is an ethical property, not just an
aesthetic one.

\WokeLang{} demonstrates that these values can be embedded in
language semantics without sacrificing the computational power that
programmers expect.  Consent gates create mandatory deliberation
points with formal safety guarantees.  Units of measure communicate
physical meaning through the type system.  Gratitude blocks make
attribution executable and auditable.  Empathy annotations surface
the social reality of development.  All of this is grounded in a
human-readable syntax designed for communication between people.

We do not claim that \WokeLang{} solves the ethical problems of
computing.  Technical solutions to ethical problems are a category
error~\citep{green2019good}.  What \WokeLang{} does is make ethical
deliberation \emph{syntactically cheap} and \emph{semantically
meaningful}, creating affordances where current languages create blind
spots.  If programming languages shape how programmers think, then a
language that makes consent, attribution, and care visible in its
syntax may---modestly, incrementally---shape programmers to think
about those things more often.

That seems worth doing.

% ============================================================================
% REFERENCES
% ============================================================================
\bibliographystyle{plainnat}

\begin{thebibliography}{99}

\bibitem[Akhawe and Felt(2013)]{akhawe2013alice}
Akhawe, D. and Felt, A.~P. (2013).
\newblock Alice in warningland: A large-scale field study of browser security
  warning effectiveness.
\newblock In \emph{Proceedings of the 22nd USENIX Security Symposium}, pages
  257--272.

\bibitem[Anderson(2008)]{anderson2008security}
Anderson, R. (2008).
\newblock \emph{Security Engineering: A Guide to Building Dependable
  Distributed Systems}.
\newblock Wiley, 2nd edition.

\bibitem[Bardzell(2010)]{bardzell2010feminist}
Bardzell, S. (2010).
\newblock Feminist {HCI}: Taking stock and outlining an agenda for design.
\newblock In \emph{Proceedings of the SIGCHI Conference on Human Factors in
  Computing Systems}, CHI '10, pages 1301--1310. ACM.

\bibitem[Bauer and Pretnar(2012)]{bauer2012programming}
Bauer, A. and Pretnar, M. (2012).
\newblock Programming with algebraic effects and handlers.
\newblock \emph{Journal of Logical and Algebraic Methods in Programming},
  84(1):108--123.

\bibitem[Bauer and Pretnar(2015)]{bauer2015programming}
Bauer, A. and Pretnar, M. (2015).
\newblock An effect system for algebraic effects and handlers.
\newblock \emph{Logical Methods in Computer Science}, 10(4).

\bibitem[Brady(2021)]{brady2021idris2}
Brady, E. (2021).
\newblock Idris~2: Quantitative type theory in practice.
\newblock In \emph{Proceedings of the 35th European Conference on
  Object-Oriented Programming (ECOOP)}.

\bibitem[Clebsch et~al.(2015)]{clebsch2015deny}
Clebsch, S., Drossopoulou, S., Blessing, S., and McNeil, A. (2015).
\newblock Deny capabilities for safe, fast actors.
\newblock In \emph{Proceedings of the 5th International Workshop on Programming
  Based on Actors, Agents, and Decentralized Control}, pages 1--12.

\bibitem[Cook(2007)]{cook2007applescript}
Cook, W.~R. (2007).
\newblock {AppleScript}.
\newblock In \emph{Proceedings of the Third ACM SIGPLAN Conference on History
  of Programming Languages}, HOPL III, pages 1--21.

\bibitem[Czaplicki(2012)]{czaplicki2012elm}
Czaplicki, E. (2012).
\newblock Elm: Concurrent {FRP} for functional {GUIs}.
\newblock Senior thesis, Harvard University.

\bibitem[{Dantec and DiSalvo}(2013)]{dantec2009infrastructuring}
{Le~Dantec}, C.~A. and DiSalvo, C. (2013).
\newblock Infrastructuring and the formation of publics in participatory
  design.
\newblock \emph{Social Studies of Science}, 43(2):241--264.

\bibitem[Dennis and Van~Horn(1966)]{dennis1966programming}
Dennis, J.~B. and Van~Horn, E.~C. (1966).
\newblock Programming semantics for multiprogrammed computations.
\newblock \emph{Communications of the ACM}, 9(3):143--155.

\bibitem[{Doctest}(2004)]{doctest2004}
{Python Software Foundation} (2004).
\newblock doctest --- Test interactive {Python} examples.
\newblock Python standard library documentation.

\bibitem[Eghbal(2020)]{eghbal2020working}
Eghbal, N. (2020).
\newblock \emph{Working in Public: The Making and Maintenance of Open Source
  Software}.
\newblock Stripe Press.

\bibitem[{European Parliament}(2016)]{gdpr2016}
{European Parliament} (2016).
\newblock Regulation ({EU}) 2016/679 ({General Data Protection Regulation}).
\newblock \emph{Official Journal of the European Union}, L119:1--88.

\bibitem[{European Parliament}(2024)]{euaiact2024}
{European Parliament} (2024).
\newblock Regulation ({EU}) 2024/1689 ({Artificial Intelligence Act}).
\newblock \emph{Official Journal of the European Union}, L2024/1689.

\bibitem[Flanagan et~al.(2005)]{flanagan2005values}
Flanagan, M., Howe, D.~C., and Nissenbaum, H. (2005).
\newblock Values at play: Design tradeoffs in socially-oriented game design.
\newblock In \emph{Proceedings of the SIGCHI Conference on Human Factors in
  Computing Systems}, CHI '05, pages 751--760.

\bibitem[Ford et~al.(2015)]{ford2015paradise}
Ford, D., Parnin, C., and Hearst, M.~A. (2015).
\newblock Paradise unplugged: Identifying themes for understanding decision
  making and affect among developers.
\newblock In \emph{Proceedings of the 8th International Workshop on Cooperative
  and Human Aspects of Software Engineering}, pages 16--19.

\bibitem[Friedman(1996)]{friedman1996value}
Friedman, B. (1996).
\newblock Value-sensitive design.
\newblock \emph{interactions}, 3(6):16--23.

\bibitem[Friedman et~al.(2013)]{friedman2013value}
Friedman, B., Kahn, P.~H., and Borning, A. (2013).
\newblock Value sensitive design and information systems.
\newblock In Doorn, N., Schuurbiers, D., van~de~Poel, I., and Gorman, M.~E.,
  editors, \emph{Early Engagement and New Technologies: Opening Up the
  Laboratory}, pages 55--95. Springer.

\bibitem[{Frink}(2002)]{frink2002}
Eliasen, A. (2002).
\newblock Frink: A practical calculating tool and programming language.
\newblock \url{https://frinklang.org/}.

\bibitem[Gabriel(1996)]{gabriel1996patterns}
Gabriel, R.~P. (1996).
\newblock \emph{Patterns of Software: Tales from the Software Community}.
\newblock Oxford University Press.

\bibitem[Gilligan(1982)]{gilligan1982different}
Gilligan, C. (1982).
\newblock \emph{In a Different Voice: Psychological Theory and Women's
  Development}.
\newblock Harvard University Press.

\bibitem[Graziotin et~al.(2014)]{graziotin2014happy}
Graziotin, D., Wang, X., and Abrahamsson, P. (2014).
\newblock Happy software developers solve problems better: Psychological
  measurements in empirical software engineering.
\newblock \emph{PeerJ}, 2:e289.

\bibitem[Green(2019)]{green2019good}
Green, B. (2019).
\newblock ``{Good}'' isn't good enough.
\newblock In \emph{AI for Social Good Workshop, NeurIPS 2019}.

\bibitem[Held(2006)]{held2006ethics}
Held, V. (2006).
\newblock \emph{The Ethics of Care: Personal, Political, and Global}.
\newblock Oxford University Press.

\bibitem[Iverson(1980)]{iverson1980notation}
Iverson, K.~E. (1980).
\newblock Notation as a tool of thought.
\newblock \emph{Communications of the ACM}, 23(8):444--465.

\bibitem[Kennedy(1997)]{kennedy1997relational}
Kennedy, A.~J. (1997).
\newblock Relational parametricity and units of measure.
\newblock In \emph{Proceedings of the 24th ACM SIGPLAN-SIGACT Symposium on
  Principles of Programming Languages}, POPL '97, pages 442--455.

\bibitem[Knuth(1984)]{knuth1984literate}
Knuth, D.~E. (1984).
\newblock Literate programming.
\newblock \emph{The Computer Journal}, 27(2):97--111.

\bibitem[Ko et~al.(2017)]{ko2017past}
Ko, A.~J., Abraham, R., Beckwith, L., Blackwell, A., Burnett, M.,~M., Erwig,
  Scaffidi, C., Lawrance, J., Lieberman, H., Myers, B., Rosson, M.~B.,
  Rothermel, G., Shaw, M., and Wiedenbeck, S. (2017).
\newblock The state of the art in end-user software engineering.
\newblock \emph{ACM Computing Surveys}, 43(3):21:1--21:44.

\bibitem[Krishnamurthi and Lerner(2019)]{krishnamurthi2019pyret}
Krishnamurthi, S. and Lerner, B.~S. (2019).
\newblock Pyret: A principled, pragmatic, and teachable programming language.
\newblock \emph{SPLASH-E}, pages 1--10.

\bibitem[Leijen(2017)]{leijen2017type}
Leijen, D. (2017).
\newblock Type directed compilation of row-typed algebraic effects.
\newblock In \emph{Proceedings of the 44th ACM SIGPLAN Symposium on Principles
  of Programming Languages}, POPL '17, pages 486--499.

\bibitem[Light(2011)]{light2011hci}
Light, A. (2011).
\newblock {HCI} as heterodoxy: Technologies of identity and the queering of
  interaction with computers.
\newblock \emph{Interacting with Computers}, 23(5):430--438.

\bibitem[Lions(1996)]{lions1996ariane}
Lions, J.-L. (1996).
\newblock {Ariane 5 Flight 501} failure: Report by the inquiry board.
\newblock Technical report, European Space Agency.

\bibitem[Miller(2006)]{miller2006robust}
Miller, M.~S. (2006).
\newblock \emph{Robust Composition: Towards a Unified Approach to Access
  Control and Concurrency Control}.
\newblock PhD thesis, Johns Hopkins University.

\bibitem[M{\"u}ller and Fritz(2015)]{muller2015stuck}
M{\"u}ller, S.~C. and Fritz, T. (2015).
\newblock Stuck and frustrated or in flow and happy: Sensing developers'
  emotions and progress.
\newblock In \emph{Proceedings of the 37th International Conference on Software
  Engineering}, pages 688--699.

\bibitem[Nelson(2006)]{nelson2006natural}
Nelson, G. (2006).
\newblock Natural language, semantic analysis and interactive fiction.
\newblock \emph{IF Theory Reader}, pages 141--188.

\bibitem[Nissenbaum(2004)]{nissenbaum2004privacy}
Nissenbaum, H. (2004).
\newblock Privacy as contextual integrity.
\newblock \emph{Washington Law Review}, 79(1):119--158.

\bibitem[Plotkin and Power(2009)]{plotkin2009handlers}
Plotkin, G. and Power, J. (2009).
\newblock Handlers of algebraic effects.
\newblock In \emph{Programming Languages and Systems}, LNCS 5502, pages
  80--94. Springer.

\bibitem[Resnick et~al.(2009)]{resnick2009scratch}
Resnick, M., Maloney, J., Monroy-Hern{\'a}ndez, A., Rusk, N., Eastmond, E.,
  Brennan, K., Millner, A., Rosenbaum, E., Silver, J., Silverman, B., and
  Kafai, Y. (2009).
\newblock Scratch: Programming for all.
\newblock \emph{Communications of the ACM}, 52(11):60--67.

\bibitem[Rigby and Bird(2013)]{rigby2013convergent}
Rigby, P.~C. and Bird, C. (2013).
\newblock Convergent contemporary software peer review practices.
\newblock In \emph{Proceedings of the 2013 9th Joint Meeting on Foundations of
  Software Engineering}, pages 202--212.

\bibitem[Sapir(1929)]{sapir1929status}
Sapir, E. (1929).
\newblock The status of linguistics as a science.
\newblock \emph{Language}, 5(4):207--214.

\bibitem[Selbst et~al.(2019)]{selbst2019fairness}
Selbst, A.~D., Boyd, D., Friedler, S.~A., Venkatasubramanian, S., and
  Vertesi, J. (2019).
\newblock Fairness and abstraction in sociotechnical systems.
\newblock In \emph{Proceedings of the Conference on Fairness, Accountability,
  and Transparency}, FAT* '19, pages 59--68.

\bibitem[Sivaramakrishnan et~al.(2021)]{sivaramakrishnan2021retrofitting}
Sivaramakrishnan, K., Dolan, S., White, L., Kelly, T., Jaffer, S., and
  Madhavapeddy, A. (2021).
\newblock Retrofitting effect handlers onto {OCaml}.
\newblock In \emph{Proceedings of the 42nd ACM SIGPLAN International Conference
  on Programming Language Design and Implementation}, PLDI '21, pages 206--221.

\bibitem[Star(1999)]{star1999ethnography}
Star, S.~L. (1999).
\newblock The ethnography of infrastructure.
\newblock \emph{American Behavioral Scientist}, 43(3):377--391.

\bibitem[Syme(2003)]{syme2003units}
Syme, D. (2003).
\newblock The {F\#} approach to units of measure.
\newblock In \emph{Trends in Functional Programming}, volume 6, pages 1--16.

\bibitem[Trist(1981)]{trist1981evolution}
Trist, E. (1981).
\newblock The evolution of socio-technical systems.
\newblock \emph{Occasional Paper}, Ontario Quality of Working Life Centre.

\bibitem[Whorf(1956)]{whorf1956language}
Whorf, B.~L. (1956).
\newblock \emph{Language, Thought, and Reality: Selected Writings of Benjamin
  Lee Whorf}.
\newblock MIT Press.

\bibitem[Winner(1980)]{winner1980artifacts}
Winner, L. (1980).
\newblock Do artifacts have politics?
\newblock \emph{Daedalus}, 109(1):121--136.

\end{thebibliography}

% ============================================================================
% APPENDIX
% ============================================================================
\appendix

\section{Complete WokeLang Example}
\label{app:complete-example}

The following program demonstrates all four novel constructs in a
realistic scenario: a user data export service that must handle
consent, units, attribution, and confidence.

\begin{lstlisting}[caption={Complete WokeLang program demonstrating all four constructs},label=lst:complete]
// SPDX-License-Identifier: MPL-2.0
// A user data export service with ethical scaffolding

module DataExport

thanks to {
  "Maria Gonzalez"   -> "Designed the GDPR-compliant export format"
  "James Okonkwo"    -> "Implemented the streaming pipeline"
  "Li Wei"           -> "Security review and consent protocol"
}

// Units for data sizes and transfer rates
remember chunk_size = 1024 measured in KB
remember max_export = 5 measured in GB
remember timeout = 30 measured in seconds

@feeling(confident=true)
to validate_export_request(user, format) ->
  remember allowed_formats = ["json", "csv", "xml"]
  match format with
    when f if f in allowed_formats -> true
    otherwise -> false

@feeling(confident=false, note="Rate limiting logic needs review")
to check_rate_limit(user) ->
  remember recent = db.count_exports(user.id, 24 measured in hours)
  return recent < 10

to export_user_data(user, format) ->
  // Validate inputs (no consent needed - read-only check)
  if not validate_export_request(user, format) {
    return error("Invalid export format")
  }

  // Consent gate: this operation accesses personal data
  only if okay "Export all personal data for {user.name}?" {
    remember data = db.query_user_data(user.id)

    // Pipeline: process the data through stages
    remember result = data
      then filter_sensitive_fields
      then format_as(format)
      then compress

    // Nested consent: sharing data externally is a separate decision
    only if okay "Send export to {user.email}?" {
      email.send(user.email, result)
      audit_log.record("export_sent", user.id, timestamp())
    }

    return result
  }

@feeling(anxious=true, note="Deletion is irreversible")
to delete_after_export(user) ->
  only if okay "PERMANENTLY delete all data for {user.name}?" {
    only if okay "This action is IRREVERSIBLE. Confirm deletion?" {
      db.delete_user(user.id)
      audit_log.record("user_deleted", user.id, timestamp())
      return "Data deleted"
    }
  }
\end{lstlisting}

\section{Consent Gate Derivation Example}
\label{app:derivation}

We show a complete typing derivation for a simple consent-gated
expression.  Let $\Gamma = \{x : \mathsf{Int}\}$, $\Sigma = \{h\}$,
and consider the expression:

\[
  e = \mathbf{consent\_handler}\; h\;
  (\mathbf{only\_if\_okay}\; \text{``delete?''}\; (x + 1))
\]

\begin{mathpar}
  \inferrule[\textsc{T-Handler}]{
    \inferrule[\textsc{T-Consent}]{
      \inferrule[\textsc{T-Add}]{
        \inferrule[\textsc{T-Var}]{x : \mathsf{Int} \in \Gamma}
          {\Gamma; \Sigma, h \vdash x : \mathsf{Int}}
        \\
        \inferrule[\textsc{T-Int}]{}
          {\Gamma; \Sigma, h \vdash 1 : \mathsf{Int}}
      }
      {\Gamma; \Sigma, h \vdash x + 1 : \mathsf{Int}}
    }
    {\Gamma; \Sigma, h \vdash
      \mathbf{only\_if\_okay}\; \text{``delete?''}\; (x + 1) :
      \mathsf{Consent}[\mathsf{Int}]}
  }
  {\Gamma; \Sigma \vdash e : \mathsf{Consent}[\mathsf{Int}]}
\end{mathpar}

\section{Comparison of Language Features}
\label{app:comparison}

\begin{table}[h]
\centering
\caption{Feature comparison with related languages}
\label{tab:feature-comparison}
\small
\begin{tabular}{@{}lcccccc@{}}
\toprule
Feature & \WokeLang{} & F\# & Inform~7 & Elm & Pyret & Pony \\
\midrule
Units of measure         & \checkmark & \checkmark &            &            &            &            \\
Natural-lang keywords    & \checkmark &            & \checkmark &            &            &            \\
Consent gates            & \checkmark &            &            &            &            &            \\
Attribution blocks       & \checkmark &            &            &            &            &            \\
Empathy annotations      & \checkmark &            &            &            &            &            \\
Helpful error messages   & \checkmark &            &            & \checkmark & \checkmark &            \\
Integrated testing       &            &            &            &            & \checkmark &            \\
Capability-based types   &            &            &            &            &            & \checkmark \\
Effect system            & Partial    &            &            &            &            & \checkmark \\
HM type inference        & \checkmark & \checkmark &            & \checkmark &            &            \\
WASM target              & Planned    &            &            & \checkmark &            &            \\
\bottomrule
\end{tabular}
\end{table}

\end{document}
